\section{Experimental design}
\label{sec:design}

\subsection{Two dataset builds and two runs}
The experiment consists of two training runs with the identical model, preprocessing and optimisation settings (Section~\ref{sec:methods}), one per dataset build (Section~\ref{sec:builds}):
\begin{itemize}
\item \textbf{Run 1} trained on build~1 (five sources, every COVID-19 image from the COVID-19 image data collection, Tuberculosis almost entirely from Shenzhen), with at most eight phase-2 epochs.
\item \textbf{Run 2} trained on build~2, which adds TBX11K and the COVID-19 Radiography Database so that COVID-19 and Tuberculosis each come from two sources, with at most fourteen phase-2 epochs and a wall-clock budget.
\end{itemize}
Both added sources were part of the original data design, with labelling rules fixed before any training, but had not finished downloading when build~1 was made. Adding them is a data intervention, not a tuning step, and it does not use the internal test split or the external cohort for any parameter. However, run~2 was started after run~1's test and external results had been seen, and those results bear directly on the question that build~2 addresses (single-source classes). The external cohort therefore cannot be considered fully untouched by the overall research process. The timeline was as follows. Run~1's external results were produced at about 09:28 on 2026-11-03. Run~2 was started at about 10:07 the same day. The rule for choosing between the runs (Section~\ref{sec:selection}) was written after run~2 had started but before any of run~2's validation, test or external results existed. We state this explicitly; a new external cohort would be needed for a strictly prospective evaluation of the selected model (Section~\ref{sec:limitations}).

\subsection{Fair comparison of the two runs}
\label{sec:selection}
Each run selected its own epoch on its own build's validation split. Because the two validation splits differ in composition, their scores are not comparable, and the test and external results may not be used to choose between runs. A selection rule was therefore fixed in advance and recorded in the repository before run~2's results existed (it was not registered with any external body), and kept unchanged in the project record.
% Source: docs/data/run_selection_protocol.md (written 2026-11-03 ~12:00); implementation lungai/evaluation/compare_runs.py -> artifacts/evaluation_reports/run_comparison.json. An image belongs to the \emph{common selection set} if (i) it is in the validation split of build~1 or build~2; (ii) it does not occur in the training or test split of either build, checked both by exact file (MD5) and by patient key (source plus patient identifier where the identity is genuine or derived from the file name, otherwise the file's SHA-256); and (iii) its label is the same in both builds where it appears in both. Montgomery images are never included.

The primary criterion is macro-F1 over the six classes on this set; the run with the higher value is selected. If the difference is below 0.01, run~1 (the model already integrated and verified in the application) is kept unless run~2's macro ROC-AUC is higher by at least 0.01. Accuracy, macro ROC-AUC, per-class recall and the same metrics restricted to the four sources present in both builds are reported alongside, because run~2 has an in-distribution advantage on the TBX11K and COVID-19 Radiography images. An exact McNemar test on the paired predictions is reported as a descriptive statistic only: with one seed per run, no claim is made that one training recipe is better in general. The run not selected is reported as a secondary experiment, labelled as such.

\subsection{Follow-up runs: seed variability and preprocessing ablation}
\label{sec:followup}
Before run~2 had been evaluated and before the two runs were compared, three further runs on the build of the selected run were fixed in the project record (Table~\ref{tab:followup}). They are trained one after another, because parallel training does not fit in 8~GB of memory.
% Source: docs/data/run_selection_protocol.md, "Follow-up runs" (fixed 2026-11-03 ~15:00).

\begin{table}[t]
\centering
\caption{Follow-up runs, fixed before the run comparison. All use the build, architecture, class weights, optimiser schedule and early-stopping patience of the selected run.}
\label{tab:followup}
\small
\begin{tabular}{@{}llll@{}}
\toprule
Run & Seed & Change from the selected recipe & Purpose \\
\midrule
Run 3 & 43 & none & Seed-to-seed variability \\
Run 4 & 44 & none & Seed-to-seed variability \\
Run 5 & \RSeed{} (as run 2) & no Gaussian pre-blur, CLAHE or low-pass filter & Preprocessing ablation \\
\bottomrule
\end{tabular}
\end{table}

Runs 3 and 4 differ from the selected run only in the seed of all random number generators (which control, for example, the initialisation of the head, the data order, augmentation and dropout). Run 5 uses the seed of the selected run but a cache of images resized and normalised without the three filters of Section~\ref{sec:preproc}, with the same manifest and the same split, so that the only difference is the filtering. Because the selected run may in principle have been ended by its wall-clock budget rather than its epoch limit, each follow-up run gets a phase-2 epoch limit equal to the number of phase-2 epochs that the selected run actually completed and a wall-clock budget well above the expected duration, so that all runs receive the same schedule; early stopping (patience 3) still applies. Each run is evaluated once on validation, test and external data.
% Ablation cache: scripts/data/make_ablation_build.py --variant no_filters.

Seed variability is reported as the mean, sample standard deviation, minimum and maximum over the selected run and runs 3 and 4. The ablation is a single run (one seed) and is read against that spread. The served model remains the one chosen by the rule in Section~\ref{sec:selection}; the follow-up runs do not replace it, because choosing the best of several seeds on test or external results would be selection on the evaluation sets.
% Aggregation: lungai/evaluation/aggregate_seeds.py -> artifacts/evaluation_reports/multi_seed_summary.json -> \RMS... and \RAbl... macros.

\subsection{Once-only evaluation}
For each run, the internal test split and the external cohort were evaluated exactly once, after the checkpoint had been fixed on validation. No threshold, preprocessing parameter or training setting was changed after those evaluations for that run. Per-class and per-source results on the test split are reported but were not acted on, with the exception discussed in the previous paragraphs (the decision to add sources for build~2, which did not use test images for any parameter).

\subsection{Analyses}
The following analyses are reported for the selected model: (a) overall internal metrics with confidence intervals; (b) per-class precision, recall, specificity, F1 and ROC-AUC, and the confusion matrix; (c) accuracy and class recall within each source of the internal test split, because the overall figure depends on the source mix and a source that is the only provider of a class can be recognised by its image style; (d) calibration (reliability diagram and ECE); (e) the external binary task, including the full table of predicted classes for Montgomery's normal and TB scans; (f) the comparison of runs on the common selection set; (g) the source-identifiability probe (Section~\ref{sec:probe}); (h) seed variability and the preprocessing ablation (Section~\ref{sec:followup}); and (i) qualitative Grad-CAM maps.

\subsection{Compute environment and reproducibility}
Training and evaluation ran on a laptop CPU with 8~GB of memory under Windows, using Python~3.12, TensorFlow~2.16 and Keras~3; TensorFlow has no GPU support on native Windows, so no GPU was used. Because of this, runs 1 and 2 are single seeds (\RSeed), seed variability is limited to the three runs of Section~\ref{sec:followup}, and every run is bounded by an epoch limit or a wall-clock budget; the run reported in Section~\ref{sec:results} trained for \REpochsTrained{} epochs and its selected checkpoint is from epoch \RSelectedEpoch. Every training run writes a metadata record containing the code version, a full configuration snapshot, the manifest SHA-256 and split fingerprint, the seed, library versions and platform, per-epoch history, the class weights and the checkpoint SHA-256. Every evaluation report repeats the checkpoint SHA-256 (\texttt{\RCheckpointShort\ldots} for the reported model, version \RModelVersion), the manifest SHA-256 and the split fingerprint. All numbers in this paper are generated from those reports into a macro file; none was typed by hand. The images themselves cannot be redistributed under the source licences; the repository contains the rules, the builder, the provenance logs and the checksums needed to rebuild the identical dataset from the official sources.

\confirm{Describe here any use of AI-assisted tools in developing the software and analysis pipeline, as required by the Elsevier policy for use of AI in the research process; see the declaration at the end of the manuscript.}
